% =====================================================================
% Apartado 2: Arquitectura (2 diapositivas, ~2 minutos, incluye SREI)
% Responsable y expositor: Apunte Pazmiño Washington German
% Claves del .bib: spotifyIncident2025, khlebnikova2024dataplatform,
%   gustavsson2023fleet, strid2023fleet, soylemez2022microservices
% Figura: figuras/arquitectura_spotify.png, exportada desde
%   diagramas/arquitectura_spotify.drawio. Sin la imagen no compila.
% SREI se describe como está implementado (cliente-servidor en capas,
%   un solo backend). No usa la nube.
% =====================================================================
\section{Arquitectura}

\begin{frame}{Cómo está organizado Spotify por dentro}
  \centering
  \includegraphics[width=\linewidth,height=0.77\textheight,keepaspectratio]{arquitectura_spotify.png}

  \vspace{0.2em}
  \raggedright
  {\scriptsize \textbf{Nube:} el tráfico entra por regiones (Europa,
  EE.\,UU., Asia-Pacífico) y el perímetro corre en contenedores que
  maneja Kubernetes \cite{spotifyIncident2025}. Elaboración propia a
  partir de
  \cite{spotifyIncident2025,khlebnikova2024dataplatform,gustavsson2023fleet,strid2023fleet}.\par}
\end{frame}

\begin{frame}{La nube en Spotify frente a nuestro SREI}
  \begin{itemize}
    \item \textbf{IaaS:} llevó su plataforma a Google Cloud con
      \emph{lift and shift} \cite{strid2023fleet}.
    \item \textbf{PaaS:} sus eventos van a colas Pub/Sub y se procesan en
      BigQuery o Dataflow \cite{khlebnikova2024dataplatform}.
  \end{itemize}

  \smallskip
  \centering
  \footnotesize
  % Celdas alineadas a la izquierda: cada una empieza con \raggedright y
  % cada fila termina con \tabularnewline (no hace falta el paquete array).
  \begin{tabular}{@{}p{0.16\linewidth}p{0.39\linewidth}p{0.39\linewidth}@{}}
    \toprule
    & \textbf{Spotify} & \textbf{SREI (nuestro PFC)} \tabularnewline
    \midrule
    Dónde corre   & \raggedright Google Cloud, por regiones
                    \cite{strid2023fleet,spotifyIncident2025}
                  & \raggedright Un solo servidor, local y para pruebas \tabularnewline
    Backend       & \raggedright Miles de componentes en microservicios
                    \cite{gustavsson2023fleet}
                  & \raggedright Una API REST en capas (Spring Boot) \tabularnewline
    Datos         & \raggedright Colas Pub/Sub y BigQuery
                    \cite{khlebnikova2024dataplatform}
                  & \raggedright PostgreSQL y carpeta \texttt{/uploads} \tabularnewline
    Si algo falla & \raggedright En la caída de 2025, Asia-Pacífico siguió
                    \cite{spotifyIncident2025}
                  & \raggedright Se cae todo: punto único de fallo \tabularnewline
    \bottomrule
  \end{tabular}

  \medskip
  \raggedright
  {\scriptsize No se trata de copiar microservicios: adoptarlos no es
  trivial y sus retos se evalúan según los requisitos del proyecto
  \cite[p.~1]{soylemez2022microservices}.\par}
\end{frame}
